\begin{abstract}
We study federated Remaining Useful Life prediction on NASA C-MAPSS turbofan data using a compact temporal convolutional model. The experimental protocol compares centralized, per-client local-only and FedAvg training with FedProx, FedAdam, an asymmetric score-aware loss and a small CNN-LSTM. It evaluates FD001--FD004 under five paired seeds, with whole-engine IID and engine-count-imbalanced splits and window-level operating-regime splits, reporting per-client and global RMSE and NASA asymmetric score. To align optimization with safety-critical maintenance risk without the gradient explosion of raw exponential objectives, we also evaluate \texttt{AsymmetricRULLoss}, a $C^1$-smooth asymmetric quadratic surrogate ($\alpha_{\text{late}}{=}2.5$, $\alpha_{\text{early}}{=}0.5$). Across a controlled five-seed comparison on FD001 and FD002, \texttt{AsymmetricRULLoss} roughly halves late-prediction rates ($40.8\%$--$63.8\%$ down to $20.5\%$--$34.2\%$) and reduces decentralized NASA scores by $37.9\%$--$48.3\%$ (FedAvg: \AsymFdOneFedNasaDelta{} on FD001, \AsymFdTwoFedNasaDelta{} on FD002), while trading off higher symmetric RMSE on FD002 due to conservative early-prediction bias. An edge-cost profiler records local training time, peak memory and model-tensor payload size; Raspberry Pi measurements were unavailable for this revision. Model parameters are exchanged without a configured TLS, differential privacy or secure aggregation mechanism, so this work makes no formal privacy claim.
\end{abstract}

\begin{IEEEkeywords}
Federated learning, predictive maintenance, remaining useful life, edge computing, C-MAPSS.
\end{IEEEkeywords}
